\color{unimain}{\subsection{Логика отключения выключателей присоединений}}
\color{black}

\begin{enumerate}

\item
Ввод и вывод логики отключения присоединений осуществляется либо посредством программного ключа \textit{«ЛО Вn»}, либо с использованием внешнего сигнала \textbf{«ЛО Вn/ ЛО: Вывод действия на Bn»}.

\item
Сигнал отключения присоединения (\textit{«ЛО Вn/ ЛО: Отключение»}) формируется при отключении СШ, а также при отключении в результате опробования присоединения. Сигнал аварийного отключения присоединения (\textit{«ЛО Вn/ ЛО: Откл. аварийно»}) дополнительно формируется при действии УРОВ на себя. Для каждого присоединения предусмотрена возможность формирования как общего, так и избирательного запрета АПВ.

\item
Общий запрет АПВ присоединения формируется по факту наличия сигнала запрет АПВ СШ. Избирательный запрет АПВ присоединения формируется по факту срабатывания ДЗО.

\item
Алгоритм работы логики отключения выключателя В1 приведен на рисунке \ref{pict_lo:Log_otkl_Q1}. Алгоритм работы логики отключения выключателей В2, В3, В4 выполнен аналогично. Параметры для настройки ЛО В1...ЛО В4 приведены в таблицах \ref{app_set:lov1}-\ref{app_set:lov4}.

\medskip
\begin{figure}[H]
\centering
\includegraphics[width=1\textwidth,height=1\textheight,keepaspectratio]{\fbpath/05.01 ДЗО/761.9901 Логика отключения присоединений/img/Log_otkl_Q1.pdf}
\captionof{figure}{Функциональная схема логики отключения В1}\label{pict_lo:Log_otkl_Q1}
\end{figure}
\medskip

\end{enumerate}